% ==============================================================================
% CAPÍTULO 5 - CONSIDERAÇÕES FINAIS
% ==============================================================================

\section{Considerações finais}
\label{sec:consideracoes}

Este trabalho apresentou o SACLA, um sistema que lê o log nativo do PostgreSQL e aponta, com
explicação, janelas de cinco minutos em que um usuário se comporta de forma atípica em relação
ao próprio histórico, usando um único algoritmo de aprendizado não supervisionado, o \iforest{}.
O problema de pesquisa perguntava se isso é possível com desempenho comparável ao de regras de
especialista e com alertas interpretáveis. A resposta é sim, com uma ressalva importante: o
modelo não supera as regras sozinho, mas as complementa.

Os objetivos específicos foram cumpridos. O registro de log foi configurado e o seu custo
medido, cerca de 10\% de latência mediana com o registro de linhas. O laboratório reproduz
quatro semanas de acessos de dez usuários sobre a base Pagila, com anomalias cujos rótulos
ficam fora do log. O \iforest{}, ajustado só na validação, obteve F1 de 68,1\% na simulação de
referência e de 66,3\% em cinco simulações confirmatórias nunca vistas; as regras, escritas
conhecendo os tipos de anomalia, ficaram perto de 80\%. A combinação dos dois detectou de 84\% a
86\% dos cenários, cerca de 11 pontos a mais que as regras sozinhas, o que foi confirmado com
teste estatístico em hipóteses registradas antes dos dados. Quando a rotina dos usuários muda
de forma legítima, as regras perdem 15,9 pontos de F1, quase todos por causa de trocas de turno,
e o \iforest{} perde 1,7.

As contribuições do trabalho são cinco. A primeira é um detector que usa só o log nativo, sem
extensões de auditoria nem alterações no SGBD, com o custo da instrumentação medido. A segunda
são os atributos relativos ao histórico de cada usuário, com um cuidado metodológico que se
mostrou decisivo: calcular os atributos das janelas de treino deixando um dia de fora, sem o que
as frações de novidade não têm efeito no modelo. A terceira são os alertas explicados pelos
atributos mais raros para o usuário e agrupados em episódios, o que tornou rápido descartar
alarmes falsos como o da rotina \textit{batch}. A quarta é um laboratório reprodutível, com
ruído legítimo, anomalias em três intensidades, um tipo inédito, mudanças de rotina e
verificações automáticas entre o que foi executado e o que foi registrado. A quinta é uma
avaliação com linha de base explícita, teste usado uma vez, réplicas, pré-registro, correção
para comparações múltiplas e relato dos resultados que não se confirmaram.

As principais limitações são o uso de dados sintéticos, desenhados pelo mesmo autor que escreveu
as regras e o catálogo de mudanças; uma base pequena e um único SGBD; o custo de latência do
registro de linhas; e a saturação do escore do \iforest{} para valores muito acima dos observados
no treino.

Como trabalhos futuros, o passo mais importante para a validade externa é aplicar o SACLA ao log
de um servidor em produção, mesmo sem rótulos, e analisar os alertas com quem conhece o sistema.
Outros caminhos são levar as regras para dentro do SACLA e alertar pela combinação dos dois
detectores, com as regras de horário reaprendidas após mudanças confirmadas na triagem;
reajustar o modelo periodicamente com uma janela móvel de histórico, o que trataria a deriva
lenta dos dados e as mudanças permanentes de rotina; avaliar o \texttt{auto\_explain} com
amostragem de comandos, medindo quanto de cobertura se preserva com menos custo; e acrescentar
atributos que expressem a distância de cada valor ao máximo do histórico do usuário, além de
avaliar variantes do algoritmo com cortes de inclinação aleatória, que reduzem os artefatos do
escore em regiões sem dados \cite{hariri2021}.
